Power electronic converters are key enabling technologies for renewable energy systems, electric vehicles, energy storage, data centres, aerospace power systems, and modern electrical grids. The rapid growth of artificial-intelligence computing is creating a particularly demanding application scenario for power conversion. Modern AI data centres \cite{11424256,11515260} require high-efficiency, high-density, highly reliable, and dynamically responsive power conversion chains to support large-scale GPU and accelerator clusters, high-current low-voltage loads, rack-level or facility-level DC distribution, energy-storage integration, and increasingly complex thermal-management systems. These requirements make converter design not only a component-level optimisation task, but also a system-level engineering problem linked to computing workload, power delivery architecture, grid interaction, manufacturability, and lifecycle reliability. In the longer term, AI-assisted engineering workflows may further close the loop between requirement definition, converter design, simulation, component selection, manufacturing planning, testing, field-data feedback, and automated redesign, enabling a more autonomous development cycle for power electronic equipment \cite{Yao2023ReAct,Schick2023Toolformer,Wu2023AutoGen}. A practical converter design therefore requires not only the selection of an appropriate circuit topology, but also coordinated decisions regarding semiconductor devices, magnetic components, capacitors, control algorithms, thermal management, electromagnetic compatibility, reliability, manufacturability, and cost \cite{10751872}.

Despite the availability of mature circuit simulation and computer-aided engineering tools, early-stage converter design remains heavily dependent on expert knowledge. Engineers often move manually between datasheets, analytical calculations, empirical design rules, loss estimation spreadsheets, magnetic design tools, thermal estimators, and simulation software. This fragmented workflow makes rapid design-space exploration difficult, especially when multiple objectives such as efficiency, power density, cost, temperature, EMI risk, and reliability must be balanced simultaneously \cite{9152082,9209087,9062664}.

Recent progress in large language models, tool-using AI agents \cite{Yao2023ReAct,Schick2023Toolformer}, and workflow-oriented multi-agent systems \cite{Wu2023AutoGen} provides an opportunity to rethink how engineering design processes can be coordinated. However, direct use of LLMs as black-box design tools is not sufficient for power electronics, where numerical accuracy, physics consistency, safety margins, and engineering traceability are essential. Therefore, a practical AI-assisted converter design system should combine LLM-based reasoning and explanation with deterministic physics models, component databases, optimisation algorithms, and simulation workflows.

This paper presents AIPE, an architecture and workflow framework for AI-assisted power electronic converter design. The core idea is to represent a converter design project as a shared structured design object and to coordinate specialised domain agents through a chief architect agent. Each domain agent performs a well-defined engineering role, while physics-based calculators, component databases, optimisation routines, simulation workflows, and human review provide quantitative checking and engineering supervision. The proposed system is intended to act as a design copilot for experienced power electronics engineers rather than as a fully autonomous replacement for engineering judgement.

The main contributions of this paper are:
\begin{enumerate}
\item A system-level AIPE architecture for AI-assisted power electronic converter design, clarifying the roles of LLM-based agents, deterministic engineering calculators, component databases, optimisation routines, simulation workflows, and human review.
\item A unified converter design object for maintaining design state, assumptions, constraints, candidate solutions, calculated metrics, risks, simulation plans, and decision history.
\item A workflow orchestration logic in which a chief architect agent coordinates specialised domain agents and manages design iteration, conflict resolution, and traceable engineering decisions.
\end{enumerate}

The remainder of this paper is organised as follows. Section II discusses the motivation and design challenges of power electronic converter design under increasingly complex AI-era power conversion requirements. Section III presents the proposed AIPE architecture and clarifies the roles of the chief architect agent, domain-specific agents, shared design state, and verification/execution support. Section IV introduces the unified converter design object used to maintain requirements, assumptions, constraints, candidate solutions, validation information, and traceable engineering decisions. Section V describes the workflow orchestration logic, including state-based iteration, conflict resolution, and feedback-driven redesign. Section VI discusses the engineering knowledge bases, deterministic calculators, rule-based checkers, and tool-chain interfaces that support physics-based screening and validation. Section VII uses a DAB isolated DC--DC converter as a running example to demonstrate structured requirement extraction, agent-level design tracing, candidate comparison, and simulation/validation planning \cite{11283728}. Section VIII discusses the role of human engineers, advantages, limitations, and the broader AIPE research series. Finally, Section IX concludes the paper and outlines future work.
